\chapter{Methodology}
\label{chap:methodology}

\section{Attention U-Net Architecture}

The segmentation backbone is the Attention U-Net of Oktay et al.~\cite{oktay2018}, adapted for 14-class abdominal CT segmentation. The architecture follows the standard encoder--decoder pattern: four downsampling blocks in the encoder reduce spatial resolution progressively while expanding channel depth, and four upsampling blocks in the decoder restore the original resolution. What distinguishes the Attention U-Net from a vanilla U-Net is the placement of a learned \emph{attention gate} at every skip connection between encoder and decoder.

Figure~\ref{fig:architecture} shows the full network diagram.

\begin{figure}[h!]
\centering
\begin{minipage}{0.90\textwidth}
  \centering
  \includegraphics[width=\textwidth]{figures/fig_unet_architecture}
  \caption{Attention U-Net architecture used in this project. The encoder (left) progressively
  downsamples through four blocks; the decoder (right) upsamples and fuses encoder features
  via attention-gated skip connections. Each block contains two $3 \times 3$ convolutions
  followed by Batch Normalisation and ReLU.}
  \label{fig:architecture}
\end{minipage}
\end{figure}

\FloatBarrier

\subsection{Encoder}

The encoder has four blocks, each comprising two $3 \times 3$ convolutional layers followed by Batch Normalisation and ReLU activations. Max-pooling with stride~2 handles downsampling between blocks. The base number of filters is $F = 32$; successive blocks double the filter count: $32 \to 64 \to 128 \to 256$.

\subsection{Bottleneck}

The bottleneck consists of two $3 \times 3$ convolutions ($256 \to 512$ channels) with a \texttt{Dropout2d} layer (drop probability $p = 0.1$) inserted between them. Spatial dropout zeroes entire feature channels rather than individual units, which regularises convolutional representations more effectively and prevents co-adaptation of adjacent filters.

\subsection{Decoder}

Each decoder block combines two inputs: an upsampled feature map from the previous decoder level and an attention-weighted encoder feature map from the corresponding skip connection. The two are concatenated along the channel dimension and passed through two $3 \times 3$ convolutional layers (BN $+$ ReLU). Upsampling uses bilinear interpolation followed by a $1 \times 1$ convolution that halves the channel count, keeping GPU memory usage manageable.

\subsection{Output Head}

A single $1 \times 1$ convolution maps the final decoder feature map to 14 channels (background plus 13 organ classes). Softmax is applied along the channel axis at inference time; during training, raw logits are passed directly to the loss function.

Table~\ref{tab:arch_params} summarises the architectural hyperparameters.

\begin{table}[h!]
\centering
\caption{Attention U-Net architectural hyperparameters.}
\label{tab:arch_params}
\begin{tabular}{ll}
\toprule
\textbf{Hyperparameter} & \textbf{Value} \\
\midrule
Input size              & $3 \times 256 \times 256$ \\
Number of encoder/decoder levels & 4 \\
Base filters ($F$)      & 32 \\
Encoder filter sequence & 32, 64, 128, 256 \\
Bottleneck filters      & 512 \\
Bottleneck dropout      & \texttt{Dropout2d}($p = 0.1$) \\
Output classes          & 14 (background + 13 organs) \\
Total parameters        & $\approx 7.85$~M \\
\bottomrule
\end{tabular}
\end{table}

\section{Attention Gate Module}

The attention gate selectively emphasises encoder features that are spatially relevant to the target organ at each decoder level. Figure~\ref{fig:attention_gate} illustrates the module.

\begin{figure}[h!]
\centering
\begin{minipage}{0.75\textwidth}
  \centering
  \includegraphics[width=\textwidth]{figures/fig_attention_gate}
  \caption{Attention gate module. The gating signal $g$ from the deeper decoder level and the
  skip feature map $x$ are projected to a common dimension, summed, passed through ReLU, and
  then projected to a scalar attention coefficient $\alpha$ via sigmoid. The output is the
  element-wise product $\alpha \odot x$.}
  \label{fig:attention_gate}
\end{minipage}
\end{figure}

\FloatBarrier

Formally, let $\mathbf{x}_l \in \mathbb{R}^{C_l \times H \times W}$ be the encoder feature map at level~$l$, and let $\mathbf{g}_l$ be the gating signal from the decoder at the coarser spatial scale. The attention coefficient is computed as:

\begin{equation}
  \alpha_l = \sigma\!\left(W_\psi \cdot \mathrm{ReLU}\!\left(
    W_x \mathbf{x}_l + W_g \mathbf{g}_l + b\right)\right)
  \label{eq:attention}
\end{equation}

where $W_x$, $W_g$, and $W_\psi$ are learned $1 \times 1$ convolutions, $b$ is a bias term, and $\sigma$ denotes the sigmoid function. The output of the gate is $\hat{\mathbf{x}}_l = \alpha_l \odot \mathbf{x}_l$, a spatially reweighted version of the encoder feature map.

Every encoder level gets its own attention gate. The intermediate dimension for the attention projections is set to $F_{\text{int}} = F_l / 2$, where $F_l$ is the channel count at level~$l$, so the network learns a distinct spatial weighting for each scale.

\section{Loss Function}

Training minimises a combined cross-entropy and Dice loss. Let $\hat{p}_{i,c}$ denote the predicted softmax probability for pixel~$i$ and class~$c$, and $y_{i,c} \in \{0, 1\}$ the one-hot ground-truth label. The combined loss is:

\begin{equation}
  \mathcal{L} = \frac{1}{2}\,\mathcal{L}_{\mathrm{CE}} + \frac{1}{2}\,\mathcal{L}_{\mathrm{Dice}}
  \label{eq:combined_loss}
\end{equation}

\noindent\textbf{Cross-Entropy term:}
\begin{equation}
  \mathcal{L}_{\mathrm{CE}} = -\frac{1}{N}\sum_{i=1}^{N}\sum_{c=0}^{13}
  y_{i,c}\,\log\hat{p}_{i,c}
\end{equation}

\noindent\textbf{Soft Dice term:} averaged over the 13 foreground classes,

\begin{equation}
  \mathcal{L}_{\mathrm{Dice}} = 1 - \frac{1}{13}\sum_{c=1}^{13}
  \frac{2\sum_i y_{i,c}\,\hat{p}_{i,c} + \epsilon}
       {\sum_i y_{i,c} + \sum_i \hat{p}_{i,c} + \epsilon}
\end{equation}

where $\epsilon = 10^{-6}$ is a smoothing constant that prevents division by zero on empty classes. The cross-entropy term stabilises early training with large gradients; the Dice term keeps supervision scale-invariant and counters class imbalance. Both terms share equal weight (0.5 each), chosen empirically based on validation Dice during preliminary runs and kept fixed across all ablation experiments.